\section{Instalación del entorno conda para MadGraph}
Tras la instalación de conda (anaconda) se procede a instalar MadGraph.
\begin{itemize}
    \item[1)] \textbf{Crear y activar el entorno conda para MadGraph}
    \begin{lstlisting}
> conda create -n MG5 -c conda-forge \
    python=3.10 \
    root \
    six \
    libxcrypt \
    compilers \
    make \
    cmake
\end{lstlisting}
    \item[2)] \textbf{Activar el entorno MG5 y verificar la versión de python}
    \begin{lstlisting}
> conda activate MG5
(MG5)> which python
(MG5)> python --version
\end{lstlisting}
    el cual debe mostrar una ruta similar a \textit{/home/usuario/anaconda3/envs/MG5/bin/python}.
    \item[3)] \textbf{Descargar MadGraph}\\
    En la página de MadGraph \url{http://madgraph.phys.ucl.ac.be/cgi-bin/Register/download.cgi} descargar la última versión estáble disponible (2.9.27).
    \item[4)] \textbf{Descomprimir el archivo \textit{MG5\_aMC\_v2.9.27.tar.gz}}\\
    Abrir la terminal en el directorio donde se encuentra el archivo descargado de MadGraph y escribir el comando
    \begin{lstlisting}
> tar -xyzf MG5_aMC_v2.9.27.tar.gz
\end{lstlisting}
    el cual debería dejar una carpeta con el nombre \textit{MG5\_aMC\_v2\_9\_27}.
\end{itemize}
